\section{Compute, runs and released artifacts}
\label{app:compute}

The cluster-based RI and intervention runs used the Tel Aviv University Slurm cluster
(partition \texttt{studentkillable}) using NVIDIA GPUs with 11--12\,GB of
memory. Each worker sharded the base OLMo-2-1124-7B model in fp16 across two
GPUs. The mature-model audit used the model without fine-tuning at revision
\texttt{7df9a82518afdecae4e8c026b27adccc8c1f0032}.
S4.3 also used Colab runtimes with A100 GPUs, including 40\,GB and 80\,GB devices;
the developmental sweep instead uses the checkpoints of Appendix~\ref{app:dev}.

Table~\ref{tab:compute} lists the production runs. Cluster intervention runs use a
frozen plan. The plan's SHA-256, the hashes of the relevant code and data files,
and the model revision are recorded in a manifest. Before measurement, a
validation gate reproduces reference candidate logits and, where applicable,
previously saved importance values within a fixed tolerance; the run aborts if
the gate fails. Output records retain per-item provenance, and interrupted runs
resume from checksummed checkpoints. Code, plans, manifests and analysis tables
are available at \url{https://github.com/Maxgolu/NLP_proj}.

\begin{figure}[h]
\centering
\resizebox{\columnwidth}{!}{%
\begin{tikzpicture}[>=stealth, font=\scriptsize, node distance=4pt,
  box/.style={draw, rounded corners=2pt, inner sep=3pt, align=center, minimum width=2.2cm, minimum height=1.6em}]
\node[box] (plan) {frozen plan\\(hashes of data, code, model)};
\node[box, right=of plan] (gate) {gate\\(reference logits, saved $I_h$)};
\node[box, right=of gate, fill=orange!25] (run) {workers\\(2 GPUs each)};
\node[box, below=8pt of run] (rec) {records with provenance\\checksummed chunks};
\node[box, left=of rec] (ana) {CPU analysis\\reproduces every table};
\node[box, left=of ana] (ver) {independent re-check\\of saved effects};
\draw[->] (plan) -- (gate); \draw[->] (gate) -- node[above, text=black!65] {pass} (run);
\draw[->] (run) -- (rec); \draw[->] (rec) -- (ana); \draw[->] (ana) -- (ver);
\draw[->, dashed] (gate.south) |- ++(0.3,-0.35) node[right, text=black!65] {fail: abort};
\end{tikzpicture}}
\caption{Cluster intervention workflow. A run starts from a frozen plan, passes a
gate that reproduces reference quantities, writes checksummed records, and its
analysis tables are recomputed on CPU from those records and re-checked
independently.}
\label{fig:pipeline}
\end{figure}

\begin{table}[h]
\centering\scriptsize
\setlength{\tabcolsep}{3pt}
\begin{tabular}{@{}llrp{2.5cm}@{}}
\toprule
Run & Job & Wall-clock & Content \\
\midrule
Stage 1, RI scan & TODO & TODO & all 1,024 heads, 534 prompts \\
Stage 1, test-only RI & 905839 & TODO & name/word controls \\
Stage 2, causal map & 888691 & 1h15m (6 GPUs) & all heads, 40 pairs; 92 heads, 178 pairs \\
Stage 2, extension & 912879 & 33m & Scope P/F, 178 pairs \\
Stage 3, characterization & TODO & $\approx$3.7h (6 GPUs) & 263,328 patched forwards \\
Stage 3, readout & 918689 & 33m & 1,920 readout records \\
Stage 4, routes and groups & TODO & TODO & TODO \\
Developmental sweep & TODO & TODO & 17 checkpoints \\
\bottomrule
\end{tabular}
\caption{Production runs. For cluster jobs, wall-clock is the Slurm elapsed time.}
\label{tab:compute}
\end{table}
